% =========================================================================
% paper_skeleton.tex — HSI-Learning 课程论文骨架（IEEEtran, 期刊版）
%
% 编译方式：
%   Overleaf：新建项目上传本文件即可（IEEEtran/booktabs/amsmath 均为内置）
%   本地：    pdflatex paper_skeleton.tex && pdflatex paper_skeleton.tex
%
% 使用约定（第 13 章）：
%   1. 每节开头的 % >>> 注释块是该节的写作任务清单，写完后删除；
%   2. 所有占位符为 \todo{} 标注（文末定义），编译时高亮可见；
%   3. 图表先引用后填图：figs/ 目录放图，\ref 全部交叉引用；
%   4. 投稿前 checklist 见文末注释块。
% =========================================================================
\documentclass[journal]{IEEEtran}

\usepackage{amsmath,amssymb,amsfonts}
\usepackage{graphicx}
\usepackage{booktabs}          % 三线表（第 12 章 12.5 节）
\usepackage{multirow}
\usepackage[colorinlistoftodos]{todonotes} % \todo 高亮占位；投稿前删除所有 \todo
\usepackage[hidelinks]{hyperref}

% ---- 占位符宏：编译时红色高亮，全部清除后才可投稿 ----------------------
\newcommand{\todo}[1]{\textcolor{red}{[TODO: #1]}}

\begin{document}

\title{Paper Title: A Concise, Specific Statement of the Contribution\todo{标题不含缩写、不含"Novel/Study of"}}

\author{First~Author, Second~Author, and Third~Author%
\thanks{Manuscript received XXX; revised XXX.（期刊模板会自动排版本信息）}%
\thanks{The authors are with XXX. (e-mail: xxx@xxx.com)}}

\maketitle

% =========================================================================
\begin{abstract}
% >>> 摘要五句模板（约 150 词，数字必须出现）：
% 背景一句 | 缺口一句 | 方法两句（含机制关键词）| 数字一句（数据集+协议+OA/AA/Kappa）| 意义一句
\todo{Abstract: background in one sentence; gap in one; method in two with
mechanism keywords; one sentence with numbers (dataset, protocol, OA/AA/$\kappa$);
closing sentence on significance.}
\end{abstract}

\begin{IEEEkeywords}
Hyperspectral image classification, deep learning, spectral--spatial
features, \todo{3--5 个关键词，按 IEEE 控制词表习惯}
\end{IEEEkeywords}

% =========================================================================
\section{Introduction}
% >>> 三段式（第 13 章 13.2 节）：问题段（为什么重要）→ 缺口段（别人做到哪、
%     缺什么，每句带 \cite，最后一句是"入场券"）→ 贡献段（3 条 bullet，
%     每条对应实验节的一张表/图，可证伪、量化）。
% >>> 写作顺序提醒：Introduction 在 Method/Experiments 之后写。

% —— 问题段 ——
Hyperspectral images (HSIs) record hundreds of contiguous spectral bands per
pixel, enabling fine-grained land-cover classification \cite{landgrebe2002}.
\todo{问题段：具体痛点（稀有类别 / 标注成本 / 计算预算）}

% —— 缺口段 ——
Deep convolutional models jointly exploit spatial and spectral information:
3-D convolutions capture joint spectral--spatial features
\cite{chen2016,zhong2018}, while 2-D counterparts are cheaper but discard
spectral detail after dimensionality reduction \cite{makantasis2015}. Hybrid
designs trade the two off explicitly \cite{roy2020hybridsn}. \todo{缺口段：
2--3 句带引用，最后一句指出本文瞄准的空档。}

% —— 贡献段 ——
The main contributions of this paper are summarized as follows:
\begin{itemize}
  \item \todo{贡献 1：模块/机制级，对应消融表的一行；}
  \item \todo{贡献 2：量化结果（数据集 + 协议 + OA/AA，较基线 +X 点）；}
  \item \todo{贡献 3：协议/代码/分析层面的可复现贡献。}
\end{itemize}

% =========================================================================
\section{Related Work}
% >>> 每段 = 第 11 章对比矩阵的一个方法家族；段内按思想线索展开；
%     每段最后一句指出该家族未解决的问题（四段汇合处 = 本文登场位置）。
%     严禁：只罗列不对比；贬低前人。

\subsection{Convolutional Approaches}
\todo{1D/2D/3D 卷积路线 \cite{hu2015,makantasis2015,chen2016}；最后一句：未解决什么。}

\subsection{Hybrid and Residual Designs}
\todo{HybridSN \cite{roy2020hybridsn}、SSRN \cite{zhong2018}；最后一句：未解决什么。}

\subsection{Attention and Beyond}
\todo{Transformer/State-space/基础模型 \cite{hong2022spectralformer}；最后一句：未解决什么。}

% =========================================================================
\section{Methodology}
% >>> 先写符号表，再给总览图（真实形状 + 枢纽高亮，参考课程图 8-1/9-1 画法），
%     核心机制放第一小节讲透；实现细节（padding 等）放实现段或脚注。
\todo{总览图：figs/overview.pdf —— 结构 + 数据流 + 创新点高亮。}

\subsection{Problem Formulation}
Let $\mathbf{X} \in \mathbb{R}^{H \times W \times B}$ denote a hyperspectral
cube and $y_i \in \{1,\dots,C\}$ the label of pixel $i$. \todo{符号系统：每个
符号第一次出现时定义；张量形状写全。}

\subsection{Proposed Module}
\todo{核心机制：伪代码见 Algorithm~\ref{alg:method}；逐层形状可列表（参考
第 8 章表 8-1，注意无 padding 的维度收缩）。}

\subsection{Implementation Details}
\todo{框架/硬件/epochs/lr/优化器——与 run\_config.json 逐项对应（第 12 章）。}

% =========================================================================
\section{Experiments}
% >>> 五件套（第 13 章 13.4 节）：Datasets and Protocols → Implementation →
%     Main Results → Ablation → Analysis。协议写进表注；每张图必须被讨论。

\subsection{Datasets and Protocols}
\todo{数据集表（尺寸/波段/类别/样本量）+ 协议表（划分/预处理/预算/种子，
对应 run\_config.json；多种子报 mean$\pm$std，第 12 章）。}

\subsection{Comparison With the State of the Art}
表~\ref{tab:main} 汇总协议 C 下的主结果（课程真实数字，替换为你自己的）。

\begin{table*}[!t]
  \centering
  \caption{Classification results under protocol C (10/10/80, seed 42) on the
  Indian Pines corrected scene. Best results are in boldface.}
  \label{tab:main}
  \begin{tabular}{lcccc}
    \toprule
    Model & Input & OA (\%) & AA (\%) & $\kappa$ \\
    \midrule
    SVM (RBF) & spectra            & 80.70 & 77.73 & 0.7795 \\
    1D CNN    & spectra            & 75.27 & 60.15 & 0.7147 \\
    3D CNN    & 9$\times$9 patch   & 94.55 & 76.53 & 0.9377 \\
    2D CNN    & 9$\times$9 patch   & 95.83 & 87.04 & 0.9525 \\
    HybridSN  & 25$\times$25 patch & 96.66 & \textbf{93.28} & 0.9619 \\
    SSRN      & 7$\times$7 patch   & \textbf{97.87} & 79.84 & \textbf{0.9756} \\
    \bottomrule
  \end{tabular}
\end{table*}

\subsection{Ablation Study}
\todo{消融表：一次一变量、mean$\pm$std（3 种子）、对照组来自同一实验矩阵
（第 12 章表；SIR 与 no-residual 消融可直接填入）。}

\begin{table}[!t]
  \centering
  \caption{Ablation under protocol D (20/10/70, mean $\pm$ std over 3 seeds).}
  \label{tab:ablation}
  \begin{tabular}{lcc}
    \toprule
    Variant & OA (\%) & AA (\%) \\
    \midrule
    Baseline (SSRN)        & \todo{xx.xx $\pm$ xx.xx} & \todo{xx.xx $\pm$ xx.xx} \\
    \quad + SIR ($\lambda{=}0.1$) & \todo{xx.xx $\pm$ xx.xx} & \todo{xx.xx $\pm$ xx.xx} \\
    \quad $-$ residual     & \todo{xx.xx $\pm$ xx.xx} & \todo{xx.xx $\pm$ xx.xx} \\
    \bottomrule
  \end{tabular}
\end{table}

\subsection{Analysis and Discussion}
\todo{混淆矩阵三步读法（第 2 章）、稀有类失败案例、注意力/特征可视化
（第 10 章）——解释性小节区分"跑分论文"与"有洞察论文"。}

% =========================================================================
\section{Conclusion}
\todo{带走什么 + 局限 + 未来工作；不重复摘要，不夸大结论。}

% =========================================================================
\bibliographystyle{IEEEtran}
% 用 Zotero/Better BibTeX 导出 references.bib（引用键规范见第 11 章 11.5 节）
\bibliography{references}
% 编译前未准备 .bib 时，可用如下占位手工条目先行排版：
% \begin{thebibliography}{10}
% \bibitem{landgrebe2002} D.~Landgrebe, ``Hyperspectral image data analysis,''
%   \emph{IEEE Signal Process. Mag.}, 2002.
% \end{thebibliography}

\end{document}

% =========================================================================
% 投稿前 checklist（逐项核对后删除本块）
%   [ ] 编译零 error、零 undefined reference（所有 \ref/\cite 不出现 "？"）
%   [ ] 所有 \todo 已删除
%   [ ] 每张图/表都在正文被 \ref 引用且被讨论
%   [ ] 主表/消融表：精度全文统一两位小数、并列最优都加粗、协议写进表注
%   [ ] 图：300 dpi、字号可读、nipy_spectral 约定（课程第 1 章）
%   [ ] 摘要含数字；贡献点与实验表一一对应
%   [ ] 协议声明：划分/种子/预处理拟合范围/训练预算（第 12 章规范）
%   [ ] requirements.lock / 代码仓库链接写入脚注（可复现性承诺）
% =========================================================================
